% UBT-AI-PROVENANCE-BEGIN
% schema: ubt-ai-provenance/v1
% tier: C_working
% ai_assistance: disclosed
% human_review: risk-based
% editorial_responsibility: Ing. David Jaroš
% policy: ../../../AI_PROVENANCE.md
% notice: Working material; exhaustive human review is not claimed.
% UBT-AI-PROVENANCE-END

\chapter{Časté otázky pro recenzenty a inženýry}

\section{Vzniká metrika projekcí biquaternionu?}
Ne. Na klasickém Lorentzově řezu je symetrizovaný součin
\[
 \frac12(E_\mu^\sharp E_\nu+E_\nu^\sharp E_\mu)
 =g_{\mu\nu}\mathbf1
\]
již centrální. Stopa, předpis reálné části, preferovaný fázový řez ani
průměrování přes vlákno nejsou součástí kanonické lokální definice.

\section{Má jedno pole s osmi reálnými složkami příliš málo složek pro metriku?}
Metrika se konstruuje ze čtyř prvních kovariantních derivací, nikoli pouze
z hodnoty $\Theta$. Jejich Lorentzovsky reálné složky tvoří tetrádu se
šestnácti prvky. Po odečtení šesti lokálních Lorentzových směrů má
diferenciál zobrazení tetrády na metriku hodnost deset.

\section{Je konexe dalším fundamentálním polem?}
V zamčené klasické architektuře nikoli. Nedegenerovaná tetráda a zadaná
torze určují metricky kompatibilní konexi. Otevřený zůstává dynamický zákon,
který torzi vybírá, a přesná kompozitní nebo pomocná realizace požadovaná
konečnou akcí.

\section{Byla Einsteinova dynamika odvozena nepodmíněně?}
Ne. Lokální kinematická reprezentovatelnost a několik podmíněných
Palatiniho výsledků jsou odvozeny, ale mikroskopický výběr a normalizace úplné
akce zůstávají otevřené. Bez tohoto odvození se stav nesmí povýšit z
podmíněného uzavření.

\section{Jak lze výpočet reprodukovat?}
Začněte od citovaného aktivního zdroje, zaznamenejte jeho předpoklady a
úroveň tvrzení, spusťte uvedený verifikátor a jeho doprovodný test pytest a
odlište algebraickou kontrolu od fyzikálního výběru. Úspěšná kontrola v systému
počítačové algebry (CAS) nebo
Leanu ověřuje pouze tvrzení, které v ní bylo zakódováno.
